\section{Proof of Withdrawer Unlinkability}
\label{sec:proof}

 \newversion{ The same  proof  \emph{strategy} described in this section applies  to the latest release.
This is because the circuits it adds ($\cktFlush$, $\cktSend$, $\cktQueueAtt$) are simply \textbf{bookkeeping} circuits that label the requests differently.
 This labeling  \textbf{does not add any user-dependent information} to the
 requests, as such they do not introduce any link to the caller.
 Furthermore the circuits are \emph{permissionless} -- they can be called by anyone, without needing proof of any secret and  \emph{anonymous} -- the
identity of the caller is not revealed at application layer.
Two changes to the circuits this section does analyse must be carried through the proof: the
marker $\refundComm(\usersk, \rid)$ becomes $\ownComm(\inIndex, \usersk)$, a keyed hash of the
same secret over the caller-chosen random index $\inIndex$ in place of $\rid$ (the source
requires $\inIndex$ to be fresh per request, which the argument must assume), and
$\cktCompleteWithdraw$ proves that marker on every verdict, so its success branch is no
longer witness-free.
}


\subsection{Proof Overview}
As we discussed in Section~\ref{def:unlink}, withdrawer unlinkability informally says the following:
a cross-chain protocol is unlinkable if, for any users $U_0,U_1$, with shielded accounts, $\candZero,\candOne$  who are  \emph{equally likely to  have called a  withdraw} at time $t$ for amount $x$,   and given the \emph{transcript of all messages exchanged up to that point\footnote{In the actual definition the adversary see the transcript of all messages exchanged before and after that particular withdraw request, but the actions after the withdraw are normalized to the event that both $\candZero$ and $\candOne$ would have performed that withdraw.}}, no one should be able to infer if the withdrawer was $\candZero$ or $\candOne$ (in polynomial time).
The same definition can be extended to every action (i.e., Swap, Supply, Redeem), except Deposit.

In this section, we provide a formal proof that  protocol $\Pi$ (Section~\ref{sec:instantiation}) achieves  withdrawer unlinkability \S~\ref{def:unlink}, assuming cryptographic assumptions stated in \S~\ref{sec:assumptions} and assuming the privacy and soundness of the Midnight shielded-coin and shielded-computation layer.

\vspace{8pt}
\noindent
\emph{Proof Intuition and Approach.}
The intuition of the proof is simple: to prove that our protocol achieves unlinkability, we need to prove that  the protocol messages generated upon execution of any action call (Withdraw, Swap, Supply, Redeem) -- except Deposit --   do not leak any information about the identity of the caller.

The approach to provide such proof is to analyze the \emph{transcript} of each call: look at every message generated by our protocol,  identify what  messages \emph{encode the identity of the caller}, and then use the security assumption of the cryptographic tools employed, to conclude that such encoding hides the identity from any polynomial time adversary.


\subsection{Intuition for the Proof of  Withdrawer-unlinkability.}
A withdraw request consists of two parts: (1) the transaction that will be executed on the public chain (2) the transactions executed on the private chain. Let us call $\shdAcct_i$, the user who is executing the $\Withdraw$  flow (i.e., $\cktStartWithdraw,\cktCompleteWithdraw$).

\vspace{5pt}
\emph{Transcript on public chain $\Lpub$ (Ethereum).}
The transaction executed  on the public chain $\req$ contains the sender $\vaultEvmAddr$, the recipient of the withdraw $\destAddr^*$ and the amount. The sender is a fixed constant $\vaultEvmAddr$, independent of the caller $\shdAcct_i$.
Hence, \emph{the request on the public chain does not contain any information about the identity of the withdrawer.}
\textbf{Note. If a user re-uses  existing  evm address for $\destAddr$ instead of a fresh one, our protocol provides no unlinkability guarantees.}


\vspace{5pt}
\emph{Transcript on private chain $\Lshd$ (Midnight).}
When a user $\shdAcct_i$ of  $\Lshd$ (Midnight) starts a withdraw of $\amt$ token,  it will first create a private payment to a public burn address $\burnAddr$.
We use Midnight libraries \MidNoteT{$\mnhl{\code{receiveShielded}}(\coinIn)$} then
\MidNoteT{$\mnhl{\code{sendImmediateShielded}}(\coinIn, \burnAddr, \amt)$}.
Under the assumption that the randomness included in $\coinIn$ is secret and unpredictable to the adversary, the  Midnight shielded-coin layer guarantee that the transcript of the above calls does not leak any information about the identity of the payer ($\shdAcct_i $).
In addition to the coins generated from the Midnight-coin layer, a withdraw call, also generates an \textbf{authentication tag} $\refundComm(\usersk_i,\rid)$, where $\usersk_i$ is the application key associate to the user who is calling withdraw.  This authentication tag is computed using a pseudorandom function $\Hash. $ By the definition (Def~\ref{def:prf}, Def~\ref{def:muprf}), the outputs of $\Hash$ do not reveal any information about the secret key $\usersk_i$ used for the evaluation.
Yet, $\Hash$ is a deterministic function, so, if we call $\refundComm(\usersk_i,\rid)$ on the same $\usersk_i$ and reuse an  $\rid$ in more than one withdraw request, then the two requests are readily linked.
This is why,  our implementation relies on the fact that there is a one-to-one mapping between a requestid $\rid$ and a request.
Hence, due to the  \textbf{pseudorandomness} of $\Hash$ and the  \textbf{uniqueness} of $\rid$, $\refundComm(\usersk_i,\rid)$ does not leak any information about the caller.
Finally, if a withdraw request fails, the caller must obtain a refund from the Midnight chain.
This is done by  \MidNoteT{$\mnhl\mintShieldedToken(\mathsf{newCoin})$}. Assuming that $\mathsf{newCoin}$ contains a fresh random nonce  ($\mintNonce$) that remains private, the privacy property of Midnight shielded-coin layer guarantees that
the recipient of the refund is protected.
In addition, every execution above is attested by a zero-knowledge proof, via the Midnight private-computation layer (Kachina).
Using the privacy guarantees of the latter, we can assume that the the transcript generated by the zk proofs, does not leak any information about the secret used in the execution of every circuit above.


\vspace{5pt}
\emph{Transcript of other actions (i.e., Swap, Supply, Redeem).}
The transcript of other actions, all follow a similar structure as Withdraw, and use the same cryptographic assumptions as withdraw. Hence we informally state that these additional transacripts do not leak any information about the caller of a withdrawal.
Formally, in the proof this is implicitly obtained in our hybrid-argument proofs.
In our hybrid experiments, in every experiment we remove a secret from the transcript, and the change affects every call (not just withdrawal).


\subsection{ Assumption on Consolidated Coins in Midnight Wallets (External to Our protocol)}
\label{sub:nofingerprint}
Our proof of security  (see hybrid $\hybrid_3$ in the formal proof), uses an assumption about user's wallet spending convention.
In the proof we assume that  \textbf{a user will always keep the coins consolidated in one coin}.

Let us explain.
At very high level, a coin on Midnight is a tuple (nonce, type, value); creating it publishes a commitment, spending it publishes a nullifier.
Spending a coin on Midnight  involves  (1) publishing the nullifiers (derived from the sender's secret key);  (2) creating a new commitment for the \emph{recipient}'s coin, and  (3) attaching a zero-knowledge proof of validity of (1) and (2).
Nullifiers do not reveal anything about the amount and the spender (to a PPT Adversary), however  the \emph{number} of nullifiers used to make a shielded payment reveal  \emph{how many coins} were spent. This information could help  identifying who is the user who is spending those coins.
So, for instance, if an adversary somehow knows that Alice has 3 coins and Bob 2 (even if he does not know the value of them), by looking at the number of nullifiers released in a  payment transaction, the adversary could determine if Alice is paying.

The above power is captured by our security game, where we give the adversary  the power to choose which actions and  transactions  honest parties perform  as part of her side channel information (See Section~\ref{sec:definition}). So, an adversary could instruct Alice to receive payments of amount 2, and then 3; and Bob to receive one payment of 5.
If Alice does not always keep her coins  consolidated in one, when Alice will spend 5, she will use two nullifers (while Bob will use only 1). The number of coins leak the fact that Alice is paying.
Therefore, in this security proof \textbf{we assume that all honest users keep their wallet consolidated}.
This means, at all time, each user keeps \textbf{one private coin }(shielded UTXO). And when it spends any amount of that coin,  she always mints a change coin.


\subsection{Formal Proof}
\begin{theorem}
\emph{If $\Hash$ satisfies Multi-user PRF Security (Definition~\ref{def:muprf}), Midnight's layer commitments satisfy Hiding (Definition~\ref{def:hiding}), Midnight's circuit outputs satisfy  zero-knowledge (Definition~\ref{def:nizk}), and Midnight's nullifiers satisfy Pseudorandomness (Definition~\ref{def:prf}),  and assuming that wallet satisfies the consolidated-coin assumption (Section~\ref{sub:nofingerprint}),
then system $\Pi$ described in Section~\ref{sec:instantiation}
achieves  withdrawer unlinkability (Definition ~\ref{def:unlink}).}
\end{theorem}


\begin{proof}

Our objective is to prove that any PPT adversary $\Adv$

$$\Pr[\unlinkgame_{\Pi,\Adv}(1^{\secpar}) = 1] \le \frac{1}{2} + \negl(\secpar)$$

with $\unlinkgame$ described in Fig.~\ref{fig:unlink}.
\newcommand{\last}{6}
We prove this statement with a standard  hybrid arguments.
We start with $\hybrid_0$ corresponding to the  real world, where
experiment  $\unlinkgame_{\Pi,\Adv}$ is executed with the real protocol $\Pi$ (described in \S~\ref{sec:instantiation}). We start by assuming that $\Adv$ wins the real-world game $\unlinkgame_{\hybrid_0,\Adv}$ with probability $p_{0}$.
Then, our goal is to reach a \emph{simulated} world where the challenger is executing a protocol that is completely simulated -- no identities/secrets are used anywhere -- and claim that $\Adv$ wins in
this simulated world with the ``same'' probability  $p_{0}$ (negligibly close). Which allows us to conclude that $p_{0}$ must be close to $1/2$ since in the simulated world the protocol messages contain no secret information, and no adversary can guess which account was picked in the game with better than $1/2$ probability.

The core of the proof is to reach the \emph{simulated world} through a sequence of hybrid worlds (or hybrid experiments),
 $\hybrid_1, \hybrid_2, \ldots$  where we run the game
 $\unlinkgame_{\hybrid_i,\Adv}$ with a modified protocol $\Pi_{\hybrid_i}$.
 In  each $\Pi_{\hybrid_i}$, the protocol is modified by removing the secrets used to compute a  cryptographic object. In each hybrid world $\hybrid_i$,
we use the cryptographic properties to claim that the probability  of $\Adv$ winning the game in this world  must be  close to  the previous world $\hybrid_{i-1}$,  up to a negligible factor.
In the last hybrid game $\hybrid_{\last}$, we reach the \emph{simulated} world, a world where  no secret is used in any computation. Since secrets are not used, in such a world, the protocol messages do not leak any information about the caller -- only the public information about the transactions do.
\emph{Since the public information do not contain the identity of the caller of a withdraw, a swap, a supply or a  redeem,}, the transcript overall does not leak any information.
In this simulated world, the probability that the adversary $\Adv$ wins the game no better than just guessing.
Hence $p_{\last}= 1/2$, and the sequence of hybrid games allows use to conclude that $p_0\approx_{\negl} p_1 \approx\ldots \approx_{\negl} 1/2$ and therefore $p_0 \approx_{\negl} 1/2$, where $\approx_{\negl}$ means equal up to a negligible factor.
We provide the formal hybrid-argument proof below.


\subsubsection{Hybrid 0}
This hybrid corresponds o the game $\unlinkgame_{\Pi,\Adv}(1^{\secpar})$ described in Figure \ref{ssec:unlink}. This is the \emph{real world}, where every oracle is implemented faithfully following protocol $\Pi$.
When participating in the real-world  experiment, $\Adv$  obtains the following transcript.


\begin{itemize} [label=$\blacksquare$]
    \item\textbf{Transcript from the oracle calls  pre-challenge}:
    The adversary has oracle access to both ledger and to cross-chain oracles  (as explained in Section~\ref{sec:definition}).

          \[(\candZero\bc \candOne\bc \tokType\bc \amt\bc \auxx^*\bc \outc^*\bc
      \mathsf{st}) \gets \Adv^{\mathcal{O}}(\pp),\]


    As such $\Adv$ can request creations of honest accounts on $\Lpub$ and $\Lshd$, and
    instruct actions on behalf of those users with the cross-chain oracles: deposit, withdraw, swap, supply, redeem.
     First,  recall that a user participating in the cross-chain system holds the following keys (these are privet to the user (challenger) and not known to the adversary).:
\begin{itemize}
    \item { \bf Application secret key:} $\usersk_i$: a private key for the cross-chain system: $\usersk_i\leftarrow\{0,1\}^{\lambda}$. This is randomly chosen upon the first Deposit request.
    \item Shielded Account  ($\callerZKey_i$, $\callerZKeySk_i$) on $\Lshd$; (\zkswap\ wallet on Midnight).
    \item Public account $\pubAddr_i$  on $\Lpub$; (used to fund the derived account).
\end{itemize}
The secret keys of the simulated users are held by the challenger who simulates the experiment.


\medskip

For convenience of the reader below we report an abridged version of the the transcripts that $\Adv$ would obtain from the interactions with the cross-chain oracles $\ODep$, $\OWd$, $\OSw$, $\OSup$, $\ORed$.

For each oracle call, we highlight:
\begin{itemize}
    \item $\Dvault$: messages posted on the contract state  on the Midnight chain.
    \item $\Dshd$: message posted on the shielded-coin layer on the Midnight chain.
    \item statements, zk proofs and relations proved by the Midnight private-circuit layer.
\end{itemize}
\begin{itemize}
    \item $\ODep(\pubAddr_i, \shdAcct_i, \tokType, \amt, \auxx, \outc)$:


\begin{TxtBox}{trn:deposit}{Transcript for the $\Deposit$ flow}
Circuit \upd{$\cktStartDeposit$} (see Circuit~\ref{ckt:deposit}) \begin{OutList}
\oitem{\Dvault} $\sigNonce \mathrel{+}= 1$; \; \upd{$\depositMap[\rid] \gets \req$; \;
  $\depositView[\rid] \gets \bigl(\userComm(\secret{\usersk_i}),\, \tokType,\, \amt\bigr)$}.
\oitem{\Dshd} $\emptyset$.
\oitem{\stmt} $(\upd{\cktStartDeposit},\; \req,\; \rid,\; \upd{\depositView[\rid]},\;
  \notif(\vaultAddr, \rid, \upd{\depositMap}),\; \comIO)$, where
  $$\req = \tokType, \amt, \vaultEvmAddr, \evmNonce, \gasEnv, \sigNonce, \keyPath$$
  and
  $\keyPath_i= \userComm(\secret{\usersk_i})$.
   \oitem{\zkproofs} \bluedot
\oitem{\Rel}  \MidNoteT{$\req.\keyPath = \userComm(\secret{\usersk_i})$
  \upd{$\;\wedge\; \depositView[\rid].\fComm = \userComm(\secret{\usersk_i})$}.}
%\oitem{\wit} $\usersk_i$.
\end{OutList}
Circuit: $\cktCompleteDeposit$ (see Circuit~\ref{ckt:completedeposit})
\begin{OutList}
\oitem{\Dvault} \upd{$\depositMap[\rid]$ removed; $\depositView[\rid]$ removed}.
  \oitem{\Dshd} \bluedot
\oitem{\stmt} $\bigl(\cktCompleteDeposit,\; \rid,\; \upd{\depositView[\rid]},\; \attSig,\;
  \tokType, \amt,\; \coinClaim\bigr)$  where
  $\coinClaim=\lnk{fn:coinCommit}\coinCommit \bigl(\freshrn{\mintNonce}, \tokType,\amt,\shdAcct_i)$.
  \oitem{\zkproofs} \bluedot
 \MidNote{\oitem{\Rel} $\userComm(\usersk_i) = \upd{\depositView[\rid].\fComm}$ \;$\wedge$\;
  $\EcdsaVf( \attSig, \mpcRespKey)$ \;$\wedge$\; and $\coinClaim=\lnk{fn:coinCommit}\coinCommit \bigl(\mintNonce, \tokType,\amt,\shdAcct)$.}
  \oitem{\wit} $(\usersk_i,\; \mintNonce,\; \shdAcct)$
\end{OutList}
\end{TxtBox}

\begin{ProcBox}
    { Note:} We expect that users will re-use  the same  application secret $\usersk_i$, and hence the same $\keyPath_i$  and $\derivedAddr_i$ for every Deposit call.
\end{ProcBox}


\item $\OWd(\shdAcct_i, \pubAddr_j, \tokType, \amt, \auxx, \outc)$
%---------------------------------------------------------------------
%  $\OWd$
%---------------------------------------------------------------------
\begin{TxtBox}{trn:withdraw}{Transcript for the $\Withdraw$ flow}
 \upd{$\cktStartWithdraw$} (see Figure~\ref{ckt:startwithdraw})
\begin{OutList}
\oitem{\Dvault} $\sigNonce \mathrel{+}= 1$; \; $\reqMap[\rid] \gets \req$; \;\\
  \upd{$\wdView[\rid] \gets \bigl(\lnk{fn:refundComm}{\refundComm(\secret{\usersk_i}, \rid)},\,
  \tokType,\, \amt\bigr)$}.
\oitem{\Dshd} \emph{Nullifiers published} (one coin spent):
  \begin{itemize}[nosep,leftmargin=1.3em,topsep=1pt]
  \item \bluedot $\coinNull\bigl(\secret{\coinFund},\; \mathsf{user}(\secret{\coinsk_i})\bigr)$ (we assume consolidated-coin):
        keyed by the caller's \textbf{secret};
  \item \bluedot $\coinNull\bigl(\freshrn{\coinIn},\; \mathsf{contract}(\vaultAddr)\bigr)$: the contract
        spending what it has just received, keyed by a \textbf{public} address.
  \end{itemize}
  \emph{Commitments published} (one coin created):
  \begin{itemize}[nosep,leftmargin=1.3em,topsep=1pt]
  \item \bluedot $\coinCommit\bigl(\freshrn{\coinIn},\; \mathsf{contract}(\vaultAddr)\bigr)$ : the wallet's caller paying the contract;
  \item \bluedot $\coinCommit\bigl(\coinBurnOut,\; \burnAddr\bigr)$, over
        $\bigl(\evolveNonce(\burnNonce),\, \mathsf{color}^{*},\, \amt\bigr)$: the contract
        paying the burn address;
  \item \bluedot $\coinCommit\bigl(\coinChange,\; \mathsf{user}(\coinpk)\bigr)$, of value
        $(V - \amt)$, recipient is the caller's public coin key.
  \end{itemize}
 \MidNote{
\oitem{\stmt} $\bigl(\upd{\cktStartWithdraw},$ $\req,$ $\rid,$ $\upd{\wdView[\rid]},$
$\Dshd,$
  $\notif(\vaultAddr, \rid, \reqMap),$ $\comIO\bigr)$, where $\req$ exposes
  $\tokType, \amt, \destAddr, \evmNonce, \sigNonce$ and $\keyPath = \keyPathLit$.}
%\WitNote{\oitem{\wit} $\bigl(\usersk,\; \burnNonce,\; \coinsk,\; \{\coinFund_i\}_i\bigr)$.}
 \MidNote{\oitem{\Rel} \upd{$\wdView[\rid].\fComm = \refundComm(\usersk, \rid)$} \;$\wedge$\;
  $\coinIn = (\burnNonce, \mathsf{color}^{*}, \amt)$ opens both
  $\coinCommit(\coinIn, \mathsf{contract}(\vaultAddr))$ and
  $\coinNull(\coinIn, \mathsf{contract}(\vaultAddr))$ \;$\wedge$\;
  each $\coinFund$ opens its published nullifier under $\coinsk$ and is unspent \;$\wedge$\;
  the color balances: $V + \amt$ in, $\amt + (V-\amt) + \amt$ out.
  }
\end{OutList}

\vspace{5pt}
$\cktCompleteWithdraw$ (See circuit~\ref{ckt:completewithdraw})
\begin{OutList}
\oitem{\Dvault} $\reqMap[\rid]$ removed; \upd{$\wdView[\rid]$ removed}.
\oitem{\Dshd} $\okbit = 1$: $\emptyset$.   $\okbit = 0$: \upd{\bluedot mint record
  $\bigl(\lnk{fn:domSep}{\domSep(\wdView[\rid].\fErc)},\, \wdView[\rid].\fAmt\bigr)$; \;
  \bluedot $\lnk{fn:coinCommit}{\coinCommit(\freshrn{\mintNonce}, \wdView[\rid].\fErc,
  \wdView[\rid].\fAmt, \shdAcct_i)}$}.
\oitem{\stmt} $\bigl(\cktCompleteWithdraw,\; \rid,\; \attSig,\; \okbit,\;
  \upd{\wdView[\rid]},\; \Dshd\bigr)$.
%\oitem{\wit} $\okbit = 1$: $\bot$. \quad $\okbit = 0$: $(\usersk, \mintNonce)$.
\oitem{\Rel} $\okbit = 1$: $\EcdsaVf(\attDigest(\rid,\evmOut), \attSig, \mpcRespKey)$. \quad
  $\okbit = 0$: the same, $\wedge$ \upd{$\refundComm(\secret{\usersk_i}, \rid) =
  \wdView[\rid].\fComm$ $\wedge$ commitment opening
  $(\freshrn{\mintNonce}, \cdot, \wdView[\rid].\fAmt, \shdAcct_i)$}.
\end{OutList}

\vspace{10pt}

\upd{$\cktRefundWithdraw$} (See circuit~\ref{ckt:refundwithdraw})
\begin{OutList}
\oitem{\Dvault}  \upd{$\reqMap[\rid]$ removed; $\wdView[\rid]$ removed}.
\oitem{\Dshd} \bluedot mint record $\bigl(\domSep(\upd{\wdView[\rid].\fErc}),\,
  \upd{\wdView[\rid].\fAmt}\bigr)$; \;
  \bluedot $\lnk{fn:coinCommit}{\coinCommit(\freshrn{\mintNonce},\tokType,\amt, \shdAcct_i)}$.
\oitem{\stmt} $\bigl(\upd{\cktRefundWithdraw},\; \rid,\; \attSig,\; \upd{\wdView[\rid]},\;
  \Dvault,\; \Dshd\bigr)$.
%\oitem{\wit} $(\usersk,\; \mintNonce)$.
\oitem{\Rel} $\refundComm(\usersk, \rid) = \upd{\wdView[\rid].\fComm}$ \;$\wedge$\;
  \upd{the relation of $\cktFailGate$} \;$\wedge$\; the commitment opens to
  $(\freshrn{\mintNonce}, \cdot, \upd{\wdView[\rid].\fAmt}, \secret{\shdAcct_i})$.
\end{OutList}
\end{TxtBox}

\item $\OSw(\shdAcct\bc \tokType\bc \amt\bc \tokTypeB\bc \amtB\bc \auxx\bc \outc)$
%---------------------------------------------------------------------
%  $\OSw$
%---------------------------------------------------------------------
\begin{TxtBox}{trn:swap}{Transcript for the $\Swap$ flow.}
 Circuit \upd{$\cktStartSwap$} (see Table~\ref{ckt:startswap})
\begin{OutList}
\oitem{\Dvault} $\sigNonce \mathrel{+}= 1$; \; $\swapMap[\rid] \gets \req$; \;
  \upd{$\swapView[\rid] \gets \bigl(\lnk{fn:refundComm}{\refundComm(\usersk, \rid)},\,
  \tokType,\, \tokTypeB,\, \amtOut,\, \amtMax\bigr)$}.
\oitem{\Dshd} as \cktref{ckt:startwithdraw}, with $\amtMax$ in place of $\amt$: \;
  \bluedot $\coinCommit(\coinIn, \vaultAddr)$, \; \bluedot $\coinNull(\coinIn, \vaultAddr)$, \;
  \bluedot $\coinCommit(\coinBurnOut, \burnAddr)$, \;
  \bluedot $\{\coinNull(\coinFund_i, \coinsk_i)\}$ and $\coinCommit(\coinChange, \secret{\srcAcct_i})$.
\oitem{\stmt} $\bigl(\upd{\cktStartSwap},$ $\req,$ $\rid,$ $\upd{\swapView[\rid]},$ $\Dshd,$
  $\notif(\vaultAddr, \rid, \swapMap),$ $\comIO\bigr)$, where $\req$ exposes
  $\tokType, \tokTypeB, \feeTier, \amtOut, \amtMax$.
\oitem{\wit} $(\usersk,\; \burnNonce,\; \coinsk)$.
\oitem{\Rel} as \cktref{ckt:startwithdraw}, over \upd{$\swapView$} and $\amtMax$.
\end{OutList}

\vspace{5pt}
Circuit $\cktCompleteSwap$ (see Table~\ref{ckt:completeswap})
\begin{OutList}
\oitem{\Dvault} $\swapMap[\rid]$ removed; \upd{$\swapView[\rid]$ removed}.
\oitem{\Dshd} \bluedot mint records $\bigl(\domSep(\tokTypeB), \amtOut\bigr)$ and
  $\bigl(\domSep(\tokType), \amtChg\bigr)$; \;
  \bluedot $\coinCommit(\freshrn{\mintNonce}, \shdAcct_i)$ and
  \upd{$\coinCommit(\freshrn{\changeNonce}, \shdAcct_i)$}.
\oitem{\stmt} $\bigl(\cktCompleteSwap,\; \rid,\; \attSig,\; \amtIn,\;
  \upd{\swapView[\rid]},\; \Dshd\bigr)$.
%\oitem{\wit} $(\usersk,\; \freshrn{\mintNonce},\; \freshrn{\changeNonce})$
\oitem{\Rel} $\refundComm(\usersk, \rid) = \upd{\swapView[\rid].\fComm}$
  \;$\wedge$\; the ECDSA statement \;$\wedge$\; \upd{$\changeNonce \neq \mintNonce$}
  \;$\wedge$\; \upd{$\amtIn \leq \amtMax$} \;$\wedge$\; the two commitments open under
  $\mintNonce$ and \upd{$\changeNonce$}.
\end{OutList}

\vspace{5pt}
\upd{Circuit $\cktRefundSwap$ (see Table~\ref{ckt:refundswap})}
 \begin{OutList}
\oitem{\Dvault} \upd{$\swapMap[\rid]$ removed; $\swapView[\rid]$ removed}.
\oitem{\Dshd} \upd{\bluedot mint record $\bigl(\domSep(\swapView[\rid].\fTokIn),\,
  \swapView[\rid].\fAmtMax\bigr)$; \;
  \bluedot $\lnk{fn:coinCommit}{\coinCommit(\freshrn{\mintNonce}, \tokType, \amtMax, \shdAcct_i)}$}.
\oitem{\stmt} \upd{$\bigl(\cktRefundSwap,\; \rid,\; \attSig,\; \swapView[\rid],\; \Dvault,\;
  \Dshd\bigr)$}.
%\oitem{\wit} $(\usersk,\; \mintNonce)$.
\oitem{\Rel} \upd{$\refundComm(\usersk, \rid) = \swapView[\rid].\fComm$ \;$\wedge$\;
  the relation of $\cktFailGate$ \;$\wedge$\; the commitment opens to
  $(\freshrn{\mintNonce}, \cdot, \swapView[\rid].\fAmtMax, \secret{\shdAcct_i})$}.
\end{OutList}
\end{TxtBox}


%%%%%%%------------


\item $\OSup(\shdAcct\bc \tokType\bc \amt\bc \auxx\bc \outc)$

\begin{TxtBox}{trn:supply}{Transcript for the $\Supply$ flow.}

 Circuit \upd{$\cktStartSupply$} (see Table~\ref{ckt:startsupply})
\begin{OutList}
\oitem{\Dvault} $\sigNonce \mathrel{+}= 1$; \; $\supplyMap[\rid] \gets \req$; \;
  \upd{$\supplyView[\rid] \gets \bigl(\lnk{fn:refundComm}{\refundComm(\secret{\usersk_i}, \rid)},\,
  \amt\bigr)$}.
\oitem{\Dshd} as \cktref{ckt:startwithdraw}, with $\colrU$ in place of $\mathsf{color}^{*}$: \;
  \bluedot $\coinCommit(\coinIn, \vaultAddr)$, \; \bluedot $\coinNull(\coinIn, \vaultAddr)$, \;
  \bluedot $\coinCommit(\coinBurnOut, \burnAddr)$, \;
  \bluedot $\{\coinNull(\coinFund_i, \coinsk_i)\}$ and $\coinCommit(\coinChange, \secret{\srcAcct_i})$.
\oitem{\stmt} $\bigl(\upd{\cktStartSupply},$ $\req,$ $\rid,$ $\upd{\supplyView[\rid]},$ $\Dshd,$
  $\notif(\vaultAddr, \rid, \supplyMap),$ $\comIO\bigr)$, where $\req$ exposes
  $\amt$, and  $\underAddr$, $\wrapAddr$, $\vaultEvmAddr$, $\evmNonce$, $\sigNonce$ and
  $\keyPath = \keyPathLit$. The notification \textbf{names $\supplyMap$}, so it alone identifies
  the request as a supply.
\oitem{\wit} $(\usersk,\; \burnNonce,\; \coinsk)$.
\oitem{\Rel} as \cktref{ckt:startwithdraw}, over \upd{$\supplyView$} and $\colrU$.
\end{OutList}

\vspace{5pt}
Circuit $\cktCompleteSupply$ (see Table~\ref{ckt:completesupply})
\begin{OutList}
\oitem{\Dvault} $\supplyMap[\rid]$ removed; \upd{$\supplyView[\rid]$ removed}.
\oitem{\Dshd} \bluedot mint record $\bigl(\lnk{fn:domSep}{\domSep(\wrapAddr)},\, \amtShares\bigr)$; \;
  \bluedot $\lnk{fn:coinCommit}{\coinCommit(\freshrn{\mintNonce}, \wrapAddr, \amtShares, \shdAcct_i)}$.
\oitem{\stmt} $\bigl(\cktCompleteSupply,\; \rid,\; \attSig,\; \amtShares,\;
  \upd{\supplyView[\rid]},\; \Dvault,\;
  \Dshd\bigr)$, where $\amtShares$ is read from $\evmOut$.
%\oitem{\wit} $(\usersk,\; \freshrn{\mintNonce})$
\oitem{\Rel} $\refundComm(\usersk, \rid) = \upd{\supplyView[\rid].\fComm}$
  \;$\wedge$\; the ECDSA statement \;$\wedge$\; the commitment opens to
  $(\freshrn{\mintNonce},\, \colrW,\, \amtShares,\, \secret{\shdAcct_i})$.
\end{OutList}

\vspace{5pt}
\upd{Circuit $\cktRefundSupply$ (see Table~\ref{ckt:refundsupply})}
\begin{OutList}
\oitem{\Dvault} \upd{$\supplyMap[\rid]$ removed; $\supplyView[\rid]$ removed}.
\oitem{\Dshd} \upd{\bluedot mint record $\bigl(\lnk{fn:domSep}{\domSep(\underAddr)},\,
  \supplyView[\rid].\fAmt\bigr)$; \;
  \bluedot $\lnk{fn:coinCommit}{\coinCommit(\freshrn{\mintNonce}, \underAddr, \amt, \shdAcct_i)}$}.
\oitem{\stmt} \upd{$\bigl(\cktRefundSupply,\; \rid,\; \attSig,\; \supplyView[\rid],\; \Dvault,\;
  \Dshd\bigr)$}.
%\oitem{\wit} $(\usersk,\; \mintNonce)$.
\oitem{\Rel} \upd{$\refundComm(\usersk, \rid) = \supplyView[\rid].\fComm$ \;$\wedge$\;
  the relation of $\cktFailGate$ \;$\wedge$\; the commitment opens to
  $(\freshrn{\mintNonce},\, \colrU,\, \supplyView[\rid].\fAmt,\, \secret{\shdAcct_i})$}.
\end{OutList}

\end{TxtBox}


\item $\ORed(\shdAcct\bc \wrapOf{\tokType}\bc \amtShares\bc \auxx\bc \outc)$
\begin{TxtBox}{trn:redeem}{Transcript for the $\Redeem$ flow.}
 Circuit \upd{$\cktStartRedeem$} (see Table~\ref{ckt:startredeem})
\begin{OutList}
\oitem{\Dvault} $\sigNonce \mathrel{+}= 1$; \; $\redeemMap[\rid] \gets \req$; \;
  \upd{$\redeemView[\rid] \gets \bigl(\lnk{fn:refundComm}{\refundComm(\secret{\usersk_i}, \rid)},\,
  \amtShares\bigr)$}.
\oitem{\Dshd} as \cktref{ckt:startsupply}, with $\colrW$ in place of $\colrU$ and $\amtShares$ in place
  of $\amt$: \;
  \bluedot $\coinCommit(\coinIn, \vaultAddr)$, \; \bluedot $\coinNull(\coinIn, \vaultAddr)$, \;
  \bluedot $\coinCommit(\coinBurnOut, \burnAddr)$, \;
  \bluedot $\{\coinNull(\coinFund_i, \coinsk_i)\}$ and $\coinCommit(\coinChange, \secret{\srcAcct_i})$.
\oitem{\stmt} $\bigl(\upd{\cktStartRedeem},$ $\req,$ $\rid,$ $\upd{\redeemView[\rid]},$ $\Dshd,$
  $\notif(\vaultAddr, \rid, \redeemMap),$ $\comIO\bigr)$, where $\req$ exposes
   $\amtShares$ and $\wrapAddr$, $\vaultEvmAddr$, $\evmNonce, \sigNonce$ and $\keyPath = \keyPathLit$. The
  notification \textbf{names $\redeemMap$}.
\oitem{\wit} $(\usersk,\; \burnNonce,\; \coinsk)$.
\oitem{\Rel} as \cktref{ckt:startwithdraw}, over \upd{$\redeemView$} and $\colrW$.
\end{OutList}

\vspace{5pt}
Circuit $\cktCompleteRedeem$ (see Table~\ref{ckt:completeredeem})
\begin{OutList}
\oitem{\Dvault} $\redeemMap[\rid]$ removed; \upd{$\redeemView[\rid]$ removed}.
\oitem{\Dshd} \bluedot mint record $\bigl(\lnk{fn:domSep}{\domSep(\underAddr)},\, \amtAssets\bigr)$; \;
  \bluedot $\lnk{fn:coinCommit}{\coinCommit(\freshrn{\mintNonce}, \underAddr, \amtAssets, \shdAcct_i)}$.
\oitem{\stmt} $\bigl(\cktCompleteRedeem,\; \rid,\; \attSig,\; \amtAssets,\;
  \upd{\redeemView[\rid]},\; \Dvault,\;
  \Dshd\bigr)$, where $\amtAssets$ is read from $\evmOut$.
%\oitem{\wit} $(\usersk,\; \freshrn{\mintNonce})$
\oitem{\Rel} $\refundComm(\usersk, \rid) = \upd{\redeemView[\rid].\fComm}$
  \;$\wedge$\; the ECDSA statement \;$\wedge$\; the commitment opens to
  $(\freshrn{\mintNonce},\, \colrU,\, \amtAssets,\, \secret{\shdAcct_i})$.
\end{OutList}

\vspace{5pt}
\upd{Circuit $\cktRefundRedeem$ (see Table~\ref{ckt:refundredeem})}
\begin{OutList}
\oitem{\Dvault} \upd{$\redeemMap[\rid]$ removed; $\redeemView[\rid]$ removed}.
\oitem{\Dshd} \upd{\bluedot mint record $\bigl(\lnk{fn:domSep}{\domSep(\wrapAddr)},\,
  \redeemView[\rid].\fShares\bigr)$; \;
  \bluedot $\lnk{fn:coinCommit}{\coinCommit(\freshrn{\mintNonce}, \wrapAddr, \amtShares, \shdAcct_i)}$}.
\oitem{\stmt} \upd{$\bigl(\cktRefundRedeem,\; \rid,\; \attSig,\; \redeemView[\rid],\; \Dvault,\;
  \Dshd\bigr)$}.
%\oitem{\wit} $(\usersk,\; \mintNonce)$.
\oitem{\Rel} \upd{$\refundComm(\usersk, \rid) = \redeemView[\rid].\fComm$ \;$\wedge$\;
  the relation of $\cktFailGate$ \;$\wedge$\; the commitment opens to
  $(\freshrn{\mintNonce},\, \colrW,\, \redeemView[\rid].\fShares,\, \secret{\shdAcct_i})$}.
\end{OutList}

\end{TxtBox}

\end{itemize}

    \item \textbf{Transcript from the challenge phase:}

    In the challenge phase, the adversary $\Adv$ selects
    two honest shielded accounts  $\candZero, \candOne$ and an amount $\amt$ to be withdrawn.

    The challenge chooses one of the two users at random ($\candB$), and  executes the withdraw call honestly using that account. Namely:
          $$\Withdraw(\candB, \freshrn{\destAddr^*}, \tokType, \amt, \auxx \,;\, \usersk_b) $$

          where \freshrn{$\destAddr^*$} is a fresh address created on the public chain.


As result of this call, the adversary will observe the following transcript


\begin{CktBox}{trn:withdrawchal}{Transcript for the $\Withdraw$ flow for user $\candB$}
   \upd{$\cktStartWithdraw$}
\begin{OutList}
\oitem{\Dvault} $\sigNonce \mathrel{+}= 1$; \; $\reqMap[\rid] \gets \req$; \;
  \upd{$\wdView[\rid] \gets \bigl(\lnk{fn:refundComm}{\refundComm(\secret{\usersk_b}, \rid)},\,
  \tokType,\, \amt\bigr)$}.
\oitem{\Dshd} \emph{Nullifiers published} (exactly one coin spent.):
  \begin{itemize}[nosep,leftmargin=1.3em,topsep=1pt]
  \item \bluedot $\coinNull\bigl(\secret{\coinFund},\; \mathsf{user}(\secret{\coinsk_b})\bigr)$,
        keyed by the caller's \textbf{secret};
  \item \bluedot $\coinNull\bigl(\freshrn{\coinIn},\; \mathsf{contract}(\vaultAddr)\bigr)$.
  \end{itemize}
  \emph{Commitments published} (one  coin created):
  \begin{itemize}[nosep,leftmargin=1.3em,topsep=1pt]
  \item \bluedot $\coinCommit\bigl(\freshrn{\coinIn},\; \mathsf{contract}(\vaultAddr)\bigr)$ - the wallet
        paying the contract; \textbf{public address again};
  \item \bluedot $\coinCommit\bigl(\coinBurnOut,\; \burnAddr\bigr)$, over
        $\bigl(\evolveNonce(\burnNonce),\, \mathsf{color}^{*},\, \amt\bigr)$ --- the contract
        paying the burn address;
  \item \bluedot $\coinCommit\bigl(\coinChange,\; \mathsf{user}(\secret{\coinpk_b})\bigr)$, of value
        $V - \amt$ - keyed by the caller's public coin key.
  \end{itemize}
 \MidNote{
  \oitem{\stmt} $\bigl(\upd{\cktStartWithdraw},$ $\req,$ $\rid,$ $\upd{\wdView[\rid]},$
 $\Dshd,$
   $\notif(\vaultAddr, \rid, \reqMap),$ $\comIO\bigr)$, where $\req$ exposes
   $\tokType, \amt, \destAddr^*, \evmNonce, \sigNonce$ and $\keyPath = \keyPathLit$.

 \oitem{\Rel} \upd{$\wdView[\rid].\fComm = \refundComm(\secret{\candB}, \rid)$} \;$\wedge$\;
   $\coinIn = (\freshrn{\burnNonce}, \mathsf{color}^{*}, \amt)$ opens both
   $\coinCommit(\freshrn{\coinIn}, \mathsf{contract}(\vaultAddr))$ and
   $\coinNull(\freshrn{\coinIn}, \mathsf{contract}(\vaultAddr))$ \;$\wedge$\;
    $\coinFund_i$ opens its published nullifier under $\coinsk$ and is unspent \;$\wedge$\;
  the color balances: $V + \amt$ in, $\amt + (V-\amt) + \amt$ out.
  }
\end{OutList}
 % \oitem{\wit} $\bigl(\secret{\usersk_b},\; \burnNonce,\; \secret{\coinsk_b},\; \{\coinFund_i\}\bigr)$.
\vspace{5pt}
$\cktCompleteWithdraw$ (See circuit~\ref{ckt:completewithdraw})
\begin{OutList}
\oitem{\Dvault} $\reqMap[\rid]$ removed; \upd{$\wdView[\rid]$ removed}.
\oitem{\Dshd} $\okbit = 1$: $\emptyset$.   $\okbit = 0$: \upd{\bluedot mint record
  $\bigl(\lnk{fn:domSep}{\domSep(\wdView[\rid].\fErc)},\, \wdView[\rid].\fAmt\bigr)$; \;
  \bluedot $\lnk{fn:coinCommit}{\coinCommit(\freshrn{\mintNonce}, \tokType, \amt, \secret{\candB})}$}.
\oitem{\stmt} $\bigl(\cktCompleteWithdraw,\; \rid,\; \attSig,\; \okbit,\;
  \upd{\wdView[\rid]},\; \Dshd\bigr)$.
$(\secret{\candB}, \mintNonce)$.
\oitem{\Rel} $\okbit = 1$: $\EcdsaVf(\attDigest(\rid,\evmOut), \attSig, \mpcRespKey)$. \quad
  $\okbit = 0$: the same, $\wedge$ \upd{$\refundComm(\secret{\candB}, \rid) =
  \wdView[\rid].\fComm$ $\wedge$ commitment opening
  $(\freshrn{\mintNonce}, \cdot, \wdView[\rid].\fAmt, \secret{\candB})$}.

%  \oitem{\wit} [private] $\okbit = 1$: $\bot$. \quad $\okbit = 0$:
\end{OutList}

\vspace{10pt}

\upd{$\cktRefundWithdraw$} (See circuit~\ref{ckt:refundwithdraw})
\begin{OutList}
\oitem{\Dvault}  \upd{$\reqMap[\rid]$ removed; $\wdView[\rid]$ removed}.
\oitem{\Dshd} \bluedot mint record $\bigl(\domSep(\upd{\wdView[\rid].\fErc}),\,
  \upd{\wdView[\rid].\fAmt}\bigr)$; \;
  \bluedot $\lnk{fn:coinCommit}{\coinCommit(\coinMint, \secret{\candB})}$.
\oitem{\stmt} $\bigl(\upd{\cktRefundWithdraw},\; \rid,\; \attSig,\; \upd{\wdView[\rid]},\;
  \Dvault,\; \Dshd\bigr)$.

\oitem{\Rel} $\refundComm(\secret{\candB}, \rid) = \upd{\wdView[\rid].\fComm}$ \;$\wedge$\;
  \upd{the relation of $\cktFailGate$} \;$\wedge$\; the commitment opens to
  $(\mintNonce, \cdot, \upd{\wdView[\rid].\fAmt}, \secret{\candB})$.

%  \oitem{\wit} [private] $(\secret{\candB},\; \mintNonce)$.
\end{OutList}
\end{CktBox}


 \item\textbf{transcript from the oracle calls post-challenge:}
    The transcript here is the same as in the pre-challenge phase under the condition specified in the definition (See Definition~\ref{sec:definition}).
\end{itemize}


Let $p_0$ be the probability that any PPT  $\Adv$ wins the game $\unlinkgame_{\hybrid_{0},\Adv}$, i.e.


$$Pr[\unlinkgame_{\hybrid_{0},\Adv} \rightarrow 1] =p_0$$

\subsubsection{\texorpdfstring{\MidNoteT{Hybrid 1  - Midnight  Shielded Computation  and Shielded Coin Layer and }}{Hybrid 1 - Midnight Shielded Computation and Shielded Coin Layer and}}

In this hybrid experiment, we use the properties offered by the Midnight Layer.
We consider a sequence of sub-hybrid, and in each one of them we replace the secrets used in a Midnight component, and use
the underlying security of the Mindight layer to claim that
our replacement is sound.
While we state precisely which assumptions we are using, we
do not provide formal reductions here, since it requires to prove the security of the Midnight components, which we assume to be secure by hypothesis assumption.


In Hybrid 1  we will use the following cryptographic propoerties that are implicit in the Midnight shielded comptuation layer.

\begin{enumerate}
    \item \textbf{ZK Proofs protect the secrets (Zero-Knowledge Property~\ref{def:nizk})}. The zk proofs associated to the statements  $\stmt$ and relation $\Rel$ satisfy the \emph{zero-knowledge property}. Technically this means that we can replace every proof  generated using the user's secrets: $\usersk_i$, $\coinsk_i$, $\mintNonce$, $\shdAcct_i$ etc, with a simulated proof  that is calculated without any secret.

    \item \textbf{Commitments protect the receivers 's of the coins (Hiding~\ref{def:hiding}).} In circuits
    $\cktCompleteDeposit$, $\cktCompleteWithdraw$, $\cktCompleteSwap$ coins are calculated as the output of
    \MidNoteT{$\opMint$},   using the recipient's identity ($\shdAcct_i$ 's public key coin $\coinpk_i$):
    \begin{align*}
    {\coinCommit(c, \mathsf{rcpt})} &= \PersistantHash\bigl(\text{cc-tag} \concat \mathsf{nonce} \concat
  \mathsf{color} \concat \mathsf{value} \concat \coinpk_i\bigr)
\end{align*}

  Assuming that the above  Midnight Commitment functions satisfy hiding when
the  $\mathsf{nonce}$ is picked at random, we can replace the recipient identity $\coinpk_i$ with the string $0$
and no adversary can detect the change.


\item\textbf{Nullifers do not leak information about the recipients (PRF~\ref{def:prf}).} Nullifiers are computed
as part of {\MidNoteT{$\opReceive$}} code,  when the user sends a coin to the contract (in $\Withdraw$ and $\Swap$ flow).
They are computed as follows:
\begin{align*}
{\coinNull(c, \mathsf{sndr})} &= \PersistantHash\bigl(\text{cn-tag} \concat \mathsf{nonce} \concat
  \mathsf{color} \concat \mathsf{value} \concat \mathsf{sndr}\bigr)
\end{align*}

Assuming that the hash function $\PersistantHash$ behave as a PRF and assuming that $\mathsf{nonce}$ remains private throughout, the nullifer protects the user $\mathsf{sndr} (\coinsk_i)$.

\item  \textbf{Additional ZSwap layer artifacts do not leak any information about the sender and receivers of a shielded transaction. }
A ZSwap transaction contains several cryptographic layers that provides the correctness and functionality guarantees for a shielded payment to execute and land correctly.
For instance, as part of the soundness of the payments,
the value involved in the coins are also committed
in Pedersen's commitments whose correctness is proved in zk via specifically designed Schnorr's zk proofs (Definition~\ref{def:hvzk}) (see~\cite{Zswap22}).
Finally, the shielded transaction contains a ciphertext, that allows the recipient to detect that the coin generated in the transaction is for them. This ciphertext is computed with a \emph{key-private}  public key encryption scheme (See Definition~\ref{def:keyprivate}. Key private means that the ciphertext does not leak any information about the public key that was used to compute it. This property is necessary for unlinkability.

\end{enumerate}


Using the above  properties,  in Hybrid 1, we make the following changes in the way the experiment  $\unlinkgame_{\Pi,\Adv}$.

\begin{enumerate}
    \item  \textbf{Hybrid $\hybrid_{1.1}$} \textbf{Replace hones tproof with  Simulated ZK proofs (use no witness)}:  In $\unlinkgame_{\hybrid_{1.1},\Adv}$ we replace all the zk proofs  created in the system (zk proofs of the circuits, and zk proofs of the SZSwap wallets),  with \emph{simulated} zk  proofs  $\zkSim_2(\crs,\td,\stmt)$  (with  $(\crs,\td) \gets \zkSim_1(1^{\secpar},\Rel)$ ) that are calculated  \emph{without any secret.}
    as per Definition~\ref{def:nizk}).
     This means that the transcript of every zk proof, \emph{does not carry any information about the identity of the caller} of the circuits executed on Midnight.

     If $\Adv$ wins the game $\unlinkgame_{\hybrid_{1.1},\Adv}$ with a probability $p_{1.1}$ that is non-negligible different than $\unlinkgame_{\hybrid_0,\Adv}$,  then $\Adv$ can be used to break the zk property of each proof.   Therefore, due the zero-knowledge Property (\S~\ref{def:nizk}), it follows:

    $$Pr[\unlinkgame_{\hybrid_{1.1},\Adv} \rightarrow 1] = p_{1.1}\approx_{\negl(\lambda)} p_0$$

    \item \textbf{Hybrid $\hybrid_{1.2}$} \textbf{Remove Recipient Identity from Coin Commitments}: in this hybrid game, the challenger replaces all the commitments representing the coins that follow the form:
\begin{align*}
{\coinCommit(c, \mathsf{rcpt})} &= \Hash\bigl(\text{cc-tag} \concat \mathsf{nonce} \concat
  \mathsf{color} \concat \mathsf{value} \concat \mathsf{rcpt}\bigr)
\end{align*}

which takes as input the identity of the recipient  $\mathsf{rcpt}$, with  the following:

\begin{align*}
{\coinCommit(\mathsf{nonce} \concat
  \mathsf{color} \concat \mathsf{value} \concat  0)}
\end{align*}
Note, in this experiment, the challenger will still maintains the private list of the raw coin-info (nonce, color, value) for each unspent coin. Such information, however,  is not inputted  in the commitment function.

     If $\Adv$ wins the game $\unlinkgame_{\hybrid_{1.2},\Adv}$ with a probability $p_{1.2}$ that is non-negligible different than $\unlinkgame_{\hybrid_{1.1},\Adv}$,  then $\Adv$ can be used to  break the hiding of the commitment scheme used by Midnight shielded-coin layer.  The reduction to hiding can go through since the proofs are simulated and since our protocol invokes the Midnight library  with secret and random nonces for each new coin minted (evolved nonces are generated by Midnight).  Therefore,  assuming that $\coinCommit$ satisfies computationally hiding property (Definition~\ref{def:hiding}, with $\csGen(1^{\secpar})$ outputing the parameters for the hash function) it follows that: $$Pr[\unlinkgame_{\hybrid_{1.2},\Adv} \rightarrow 1]=  p_{1.2}\approx p_{1.1}$$


 \item \textbf{Hybrid $\hybrid_{1.3}$} Nullifier are computed without any sender's information ($\coinsk_i$).

 Namely, in   $\hybrid_{1.1}$ the challenge replaces this computation:

\begin{align*}
{\coinNull(c, \mathsf{sndr})} &= \PersistantHash\bigl(\text{cn-tag} \concat \mathsf{nonce} \concat
  \mathsf{color} \concat \mathsf{value} \concat \mathsf{sndr}\bigr)
\end{align*}

with  this:
\begin{align*}
{\coinNull(c, \mathsf{sndr})} &= \PersistantHash\bigl(\text{cn-tag} \concat \mathsf{nonce}  \concat \mathsf{value}\concat  0 \bigr)
\end{align*}

This computation is exactly as the coin commitment with the only difference in the tag   \textsf{cn-tag} (vs \textsf{cc-tag}). Here $\mathsf{nonce}$  is the randomness used to compute the coin commitment is some earlier invocation of $\opMint$.

 If $\Adv$ wins the game $\unlinkgame_{\hybrid_{1.3},\Adv}$ with a probability $p_{1.3}$ that is non-negligible different than $\unlinkgame_{\hybrid_{1.2},\Adv}$,  then $\Adv$ can be used to break the security of Midnight's nullifers.
 This is part of the Midnight design and we defer the reader to their proofs of security.

Therefore, $$Pr[\unlinkgame_{\hybrid_{1.3},\Adv}  \rightarrow 1]=  p_{1.3}\approx_{\negl(\lambda)}  p_{1.2}$$
\end{enumerate}


 \item \textbf{Hybrid $\hybrid_{1.4}$ \MidNote{ZSwap Wallet Layer (Pedersen Commitments and Key-Private Public key encryption.)} }
A shielded payment involves the use of other cryptographic objects  such as Pedersen's Commitments (Def.~\ref{def:commit}), Key-Private Public key Ciphertexts (Def.~\ref{def:keyprivate}), and Schnorr's proofs (Def.~\ref{def:hvzk}), as described in~\cite{Zswap22}.
Pedersen's Commitments  and Schnorr's proofs are used as part of the Midnight shielded-coin layer for the soundness of the proof, whereas Key-Private Public key Ciphertexts are used to
guarantee discoverability of the coin  to the receiver.
In this hybrid experiments, we  remove the secrets used in any of such objects. We can do so, since all zk proofs are simulated.
Our code does not interfere with this layer, neither provides the randomness used in the wallet itself, so we inherit the security properties guaranteed by the Midnight shielded coin layer.
The  resulting transcript should  be indistinguishable from the transcript generated in  $\hybrid_{1.3}$,  due to the zk property of Schnorr's zk proofs and  the Key-Privacy of the  Public key Ciphertexts.


Therefore, $$Pr[\unlinkgame_{\hybrid_{1.4},\Adv}  \rightarrow 1] = p_{1.3}\approx_{\negl(\lambda)} p_{1.3}$$


\subsubsection{Hybrid 2. Replace \texorpdfstring{$\Hash(\usersk_i, \cdot)$}{HashK(userSk\_i, .)} with a Random Function \texorpdfstring{$\RandFun_i$}{RF\_i} (uses no secret)}
In hybrid $\hybrid_2$ we change the calculation of $\userComm$ and $\refundComm$ that are directly  dependent on the user secret $\usersk_i$.
We replace the function $\Hash_{\usersk_i}$ keyed with the secret key of the user $\usersk_i$, with a random function $\RandFun_i$. Precisely,  the following computations (See section~\ref{ssec:userskeys}):
\begin{align*}
{\userComm(\usersk)_i} &= \Hash\bigl(\text{``vault:user:''}\bc \usersk_i\bigr)
\\
{\refundComm(\usersk_i, \rid)} &= \Hash\bigl(\text{``vault:refund:''}\bc \usersk_i\bc \rid\bigr)
  &&\text{\small the gate on refunds and swap settlement}
\end{align*}

are replaced with:


\begin{align*}
 {\userComm^{\RandFun_i}} &= \RandFun_i\bigl(\text{``vault:user:''}\bigr)\\
{\refundComm^{\RandFun_i}(\rid)} &= \RandFun_i\bigl(\text{``vault:refund:''}\bc  \rid\bigr)
 \end{align*}


The only difference between the hybrid  $\unlinkgame_{\hybrid_{1.4},\Adv}$  and $\unlinkgame_{\hybrid_{2},\Adv}$
is in the use of truly random functions
$\RandFun_1(\cdot)$, $\RandFun_2(\cdot), \dots, \RandFun_{\numusers}(\cdot)$ instead of
$\Hash_{\usersk_1}(\cdot)$, $\Hash_{\usersk_2}(\cdot)$, $\ldots,\Hash_{\usersk_{\numusers}}(\cdot)$.

Assume that $\Adv$ wins the game $\unlinkgame_{\hybrid_{2},\Adv}$  (when $\Hash$  was used ) with a probability $p_{2}$,  that is non-negligible different from $\unlinkgame_{\hybrid_{1.4},\Adv}$ (where independent random functions $\RandFun_i$ are used), which is $p_{1.4}$.
Then we can use $\Adv$ to build a  distinguisher $\distprf$ that violates the multi-user PRF security of $\Hash$  (Definition~\ref{def:muprf}) with  the same probability.

\begin{tcolorbox}[
  breakable,
  colback=boxbg,
  colframe=boxrule,
  arc=2pt,                 % Corner radius
  boxrule=0.8pt,           % Border thickness
  top=8pt, bottom=8pt,     % Inner padding
  left=10pt, right=10pt
]
\underline{\textbf{Reduction $\distprf^{\oracle(\cdot,\cdot)}$}}\\
Activate adversary $\Adv$ and simulated game $\unlinkgame_{\hybrid_{1.4},\Adv}$ to it as follows:
\begin{itemize}
    \item Simulate the execution of every oracle using the same approach as in Experiment $\hybrid_{1.4}$ with the following exception:
        \begin{enumerate}
            \item For every new user $U_i$ created for the cross-chain system, instead of picking the application secret key $\usersk_i$, the challenger calls the oracle on input $i$: $\oracle(i,\cdot)$
            \item For every deposit flow for user $U_i$ compute function $\userComm$ as:
            \begin{align*}
                & \userComm^{i} = \oracle(i, \text{``vault:user:''})\\
            \end{align*}
            \item For every withdraw,  Swap,  Supply, Redeem flow for user $U_i$ compute function as
            \begin{align*}
                &\refundComm^{i}(\rid) = \oracle(i,\text{``vault:refund:''}|| \rid)\\
            \end{align*}
        \end{enumerate}
    \item Challenge \[(\candZero\bc \candOne\bc \tokType\bc \amt\bc \auxx^*\bc \outc^*\bc \mathsf{st}) \gets \Adv^{\mathcal{O}}(\pp),\]
      where $\candZero, \candOne$ are honest shielded accounts.

    \textbf{Challenge}
    \begin{enumerate}
        \item If $\candZero$ or $\candOne$ is not honest, output $0$;
        \item if $\bal[\candZero][\tokType] < \amt$ or $\bal[\candOne][\tokType] < \amt$, output $0$;
    %    \item if $\Leak(\candZero, \tokType, \amt) \neq \Leak(\candOne, \tokType, \amt)$, output $0$.
        \item otherwise, pick a random bit $b$ and execute:
        \[ \Withdraw(\candB, \destAddr^*, \tokType, \amt, \auxx \,;\, \privIn_b) \]
        for a fresh public-chain address $\destAddr^*$, replacing the function $\Hash$ for $\candB$ with the evaluation of its correspondent oracle.
\end{enumerate}

\item Post-Challenge: simulate all the cross-chain oracles as above.
\item Output: upon receiving a bit $b'$ from $\Adv$.
\item If $b'=b$ output $1$, else output $0$.
\end{itemize}
\end{tcolorbox}


\vspace{5pt}
\noindent
\textbf{Analysis of $\distprf$'s probability of success:}
\begin{itemize}
    \item Case $\oracle(i,\cdot)= \RandFun_i(\cdot)$.

    If the Oracle is implemented as a collection of truly random functions, then  $\distprf$ is simulating the experiment $\unlinkgame_{\hybrid_{2},\Adv}$.
    Hence,


$$Pr[\distprf^{\RandFun_i(\cdot)}\rightarrow 1 ] = Pr[\unlinkgame_{\hybrid_{2},\Adv}\rightarrow 1 ]= p_{2}$$


     \item Case $\oracle(i,\cdot)= \Hash(sk_i,\cdot)$
      If the Oracle is implemented as the function $\Hash$ evaluated on  independently sampled secret keys $\usersk_i$, then the transcript generated by $\distprf$ is identical to $\unlinkgame_{\hybrid_{1.4},\Adv}$. Hence,


     $$Pr[\distprf^{\Hash(sk_i, \cdot)}\rightarrow 1] = Pr[\unlinkgame_{\hybrid_{1.4},\Adv}\rightarrow 1 ]= p_{1.4}$$

\end{itemize}


Hence:

$$ \left| Pr[\distprf^{\RandFun_i(\cdot)}\rightarrow 1] - Pr[\distprf^{\Hash(sk_i, \cdot)}\rightarrow 1] \right|  = p_{2}- p_{1.4}  $$

Since we assume that $\Hash$ satisfies the property defined in Definition~\ref{def:muprf}, it follows that $p_{2}- p_{1.4}$ must  be negligible.


\vspace{10pt}
We have removed the secret key of the caller from the authentication tags $\userComm$ and $\refundComm$.
In this hybrid world, the application secret key  $\usersk_i$ is not used in \emph{any computation}.


Note, however, that even  if we are using a random function $\RandFun_i$, the output  is still deterministic -- i.e., when queried with the same input, $\RandFun_i$ will return the same output.
So, for instance, since in Deposit,  $\userComm_i$ is computed on the same string ``vaul:user' every time, the adversary can easily detect that the same user is performing various deposits. This is by design: we are not protecting the linkability of the deposit.

We are interested in guaranteed that  withdrawals/swaps/supply/redeem  are unlinkable to each other and unlinkable to the deposits. Recall that for these action, the user authenticate with $\refundComm_i$.
And recall that:

$${\refundComm^{\RandFun_i}(\rid)} = \RandFun_i\bigl(\text{``vault:refund:''}\bc  \rid\bigr)$$


 where $\rid$ is unique per request (and $\rid = \Kec(\req)$).

This means that every $\refundComm_j$ is  different from any other $\refundComm_j$ for any $i,j$, and carries no information about the recipient.
Finally the request itself: $\req$
\begin{enumerate}
    \item $\req$ does not contain information about the identity of the withdrawer/swapper/supplier/redeemer.
    \item each $\req$ is assumed to be  unique due to the correctness of the underlying  EVM nounce.
\end{enumerate}


Putting everything together, in $\hybrid$ 2 we have:

\begin{enumerate}[nosep,leftmargin=1.7em,topsep=3pt]
\item[(i)]  no identity is leaked from the cryptographic components. In the previous hybrid we removed the uses of the secret inputs from the coin and circuit computation, and in this hybrid we removed the identity $\usersk_i$.
So, in hybrid $\hybrid_2$ we are left with  no secret and no account identity is used in producing any component of the view: proofs are simulated, coin commitments carry a dummy recipient, nullifiers carry a  dummy sender, and $\userComm$ and $\refundComm$ are evaluations of a random function;

\item[(ii)] no information identity is  leaked from the withdraw/swap/supply/redeem request $\req$ (this is by design);
\item[(ii)] no information is leaked in the spending  due to our consolidate-coin assumption (\S~\ref{sub:nofingerprint}).
\end{enumerate}


This means that in  the world $\hybrid$ 2, the protocol messages    are computed without any user-secret,  the public transcript contain no information that identify the user who is chosen in the challenge phase by the challenger to execute $\Withdraw$
Consequently, in this world, no adversary can guess which user was chosen  with a probability better than  $\frac{1}{2}$.
Hence:  $$\Pr[\unlinkgame_{\hybrid_2,\Adv} \rightarrow 1] = \frac{1}{2}$$


and due to the sequence of indistinguishable hybrid arguments,  we conclude:

$$\Pr[\unlinkgame_{\hybrid_0,\Adv} \rightarrow 1] = \frac{1}{2}+\negl(\secpar)$$


\qed
\end{proof}


\subsection{Swap/Supply/ Redeem}
\label{sub:swapsupplyredeemproof}


The proof of Withdraw Unlinkability develops a sequence of arguments
that removes  \emph{all user-dependent secrets $\usersk_i$ and $\shdAcct_i$ from every computation.}
The final hybrid experiment of the proof is a \emph{simulated} world where no secrets are used (in any of the procedures).
We notice that that Swap, Supply and Redeem follow the same structure of the Withdraw, and the same sequence of arguments apply.
In the following we examine each call and explain why the same proof argument applies.


\medskip
\noindent\textbf{$\Swap$.}  We notice that the circuits  $\cktStartSwap$/ $\cktCompleteSwap$
 follows the same structure  of  $\cktStartWithdraw$/ $\cktCompleteWithdraw$.
Specifically, $\cktStartSwap$ is the same as $\cktWithdraw$ with $\swapMap$ and
$\swapMarkerMap$ in place of $\reqMap$ and $\wdMarkerMap$. The logic is also
almost identical: a coin is surrendered and burned in the same two moves, the same authentication tag $\refundComm(\usersk_b, \rid)$  is computed.
The only difference is that in swap, the transaction on the public chain involves only the vault account and concerns swapping coin \emph{types}, whereas in the withdraw request, the transaction must have an explicit evm destination address $\destAddr$.
The settlement call   $\cktCompleteSwap$
differs in two ways from $\cktCompleteWithdraw$, however, neither of them depends on the user's identity.
$\cktCompleteWithdraw$  mints on the \emph{success} path rather than the failure path, and it mints
\emph{two} coins rather than one.
Minting more coins does not change our arguments in  $\hybrid_{1.1}$, since in that hybrid we replace  every proof in the transcript and $\hybrid_{1.2}$ over \emph{every coin commitment}, so a second commitment is replaced by the same hybrid and the same reduction.
The amounts released and swapped, are in the public state, just like the amount withdrawn.
 Every step of $\hybrid_0$ through $\hybrid_2$ therefore applies verbatim, with
$\lendunlinkgame$ in place of $\unlinkgame$, and Definition~\ref{def:lendunlink}.

\medskip
\noindent\textbf{$\Supply$.}  $\Supply$ is a mirror of $\Swap$ with the following differences:
(1) in $\Supply$ the token types and color are hardwired in the call (whereas in $\Swap$ they are an input chosen by the caller);  since types and color are public in all the calls, this change does not affect the analysis.
(2) $\Supply$ returns $\amtShares$ shares instead of an swapped amount.
$\amtShares$ is not chosen by the challenger, it is determined by $\amt$ and by the market's exchange rate at the
executing block, and the rate is a property of $\Lpub$, global to all suppliers, not of the shielded
account that paid.
Both differences do not affect the creation of the cryptographic objects and do not affect the analysis.

\medskip
\noindent\textbf{$\Redeem$.}  $\Redeem$ is  identical to $\Supply$ but with the color swapped, with the  released $\amtAssets$ being a function of $\amtShares$ and the same global rate.


\remove{
We formalize this in the following Lemma~\ref{lem:transfer}
\begin{lemma}[The chain of hybrid is  ``Procedure Independent'']
\label{lem:transfer}
Let $\mathsf{P}$ in ($\Swap, \Supply, \Redeem$)  be a procedure of $\Pi$ , and let
$\mathsf{Exp}_{\mathsf{P}}$ be the experiment of Fig.~\ref{fig:unlinkGeneral} with $\Withdraw$ replaced by $\mathsf{P}$. If:
\begin{enumerate}[leftmargin=2.2em,itemsep=2pt,topsep=3pt]
\item[\textbf{(C1)}] the request call writes $\keyPath = \keyPathLit$, so that every field of the stored record $\req$ is a function of the transaction arguments and of public state, and \emph{no field of it depends on the caller};
\item[\textbf{(C2)}] the only value the call publishes that depends on the caller's long-term secret is $\refundComm(\usersk, \rid)$, stored in a marker map whose identity is determined by the product and not by the caller;
\item[\textbf{(C3)}] every value released by the settlement call is a function of $\req$ and of the  public state of $\Lpub$, and the coin-layer components of the view consist of coin commitments  and nullifiers in numbers fixed  (by assumption on the spending fingerprint (Section ~\ref{sub:nofingerprint}))
\end{enumerate}
Then $\mathsf{P}$ satisfies unlinkability under the same assumptions as Theorem~1, with the same chain of hybrids and the same reductions.
\end{lemma}

\begin{proof}[Proof sketch]
In the proof \S\ref{sec:proof} each hybrid  is a substitution applied to a class of published objects and does not depend on which procedure produce them.
For instance,  $\hybrid_{1.1}$ replaces every proof by a simulated
one -- across all procedures that invoke zk proofs--;  $\hybrid_{1.2}$ replaces the recipient slot of \emph{every coin commitment} with $0$, $\hybrid_{1.3}$ replaces the sender
slot of every nullifier with $0$, $\hybrid_{1.4}$ replaces the identity from every ZSwap wallet objects (Pedersen's commitment, Schorr's proof  and Ciphertexts),
and finally  $\hybrid_2$  replaces every evaluation of $\Hash_{\usersk_i}$ with a random function $\RandFun_i$.
The only object that are specific to each procedure, are not the public parameters of the request and the public consequence of each request and do not depend on $b$. Conditions
(C1)--(C3) are exactly this claim. (C1) removes the request record, (C2) confines the caller's secret
to a single value of the form the hybrids already retire, and (C3) removes the settlement release and
fixes the coin-layer cardinalities. Clause (ii) of $\hybrid_3$ is discharged by  the no-spending fingerprint assumption \S~\ref{sub:nofingerprint}.
\end{proof}
}
